\section{Descrizione Informale del Contesto e degli Obiettivi}

La nostra applicazione mira a offrire una piattaforma web per il monitoraggio remoto e l'assistenza domiciliare di pazienti anziani. L'invecchiamento della popolazione richiede strumenti digitali accessibili, ma molte soluzioni esistenti risultano troppo complesse per l'utenza anziana oppure troppo frammentate per essere davvero utili agli operatori sanitari.

Il progetto \textbf{A.L.I.C.E.} nasce per colmare questo divario. Il sistema permette all'operatore sanitario di gestire i propri pazienti, assegnare esercizi fisici, configurare farmaci e monitorare umore, aderenza farmacologica e completamento degli esercizi. Il paziente, invece, utilizza un'interfaccia semplificata con card grandi, testi brevi e percorsi lineari.

\subsection{Doppia vista del sistema}

La piattaforma e' organizzata in due aree distinte:

\begin{itemize}
    \item \textbf{Area Operatore Sanitario:} comprende login, registrazione, dashboard con elenco pazienti, filtri per stato clinico, ricerca per nome, schede riepilogative e pagina di dettaglio del paziente. Nel dettaglio l'operatore puo' consultare metriche, storico umore, storico farmaci, storico esercizi, modificare i dati anagrafici, aggiungere o rimuovere farmaci ed esercizi.
    \item \textbf{Area Paziente Anziano:} comprende login, registrazione, home semplificata e tre funzionalita' principali: registrazione dell'umore quotidiano, esecuzione di esercizi fisici guidati passo-passo e conferma dell'assunzione dei farmaci.
\end{itemize}

\subsection{Funzionalita' principali}

Le funzionalita' effettivamente implementate nel progetto sono:

\begin{itemize}
    \item \textbf{Gestione utenza:} registrazione e login separati per operatore e paziente, con controllo del ruolo tramite Supabase Auth e tabella \texttt{profiles}.
    \item \textbf{Gestione pazienti:} ogni paziente e' associato a un operatore. L'operatore puo' visualizzare la lista, filtrarla, cercare per nome e accedere al dettaglio.
    \item \textbf{Gestione farmaci:} l'operatore inserisce nome, dosaggio e orario; il sistema calcola automaticamente la fascia oraria (\texttt{mattina}, \texttt{pranzo}, \texttt{sera}). Il paziente conferma l'assunzione giornaliera.
    \item \textbf{Gestione esercizi:} l'operatore crea esercizi fisici con durata, frequenza e passaggi guidati. Il paziente li esegue step-by-step e registra il completamento.
    \item \textbf{Monitoraggio umore:} il paziente registra ogni giorno il proprio umore scegliendo tra \texttt{triste}, \texttt{normale} e \texttt{felice}. L'operatore visualizza lo storico settimanale.
    \item \textbf{Suggerimento AI per l'operatore:} nella pagina di dettaglio paziente e' presente una route API che invia a Groq dati sintetici sul paziente e riceve una frase operativa di supporto.
\end{itemize}

\subsection{Obiettivi del progetto}

Gli obiettivi principali sono:

\begin{enumerate}
    \item \textbf{Semplificare l'interazione del paziente anziano}, riducendo la navigazione a pochi passaggi e usando componenti visivi chiari.
    \item \textbf{Centralizzare le informazioni utili all'operatore}, evitando che umore, farmaci ed esercizi siano gestiti in strumenti separati.
    \item \textbf{Misurare l'aderenza quotidiana}, tramite metriche percentuali su farmaci presi ed esercizi completati.
    \item \textbf{Migliorare la continuita' assistenziale}, permettendo all'operatore di controllare l'andamento del paziente anche tra una visita e l'altra.
    \item \textbf{Mantenere una struttura tecnica semplice}, basata su Next.js, React, Server Actions e Supabase, cosi' da rendere il progetto comprensibile e manutenibile.
\end{enumerate}

\newpage
